\documentclass{article}
\usepackage{icmlko}
\usepackage{graphicx}
\usepackage{booktabs}
\usepackage{amsmath}

\title{자를 세 번 고친 뒤에 처음 재 봤다 \\
  {\large 그리고 그 자는 처치를 아무렇게나 섞어도 같은 값을 냈다}}

\begin{document}
\maketitle

\begin{abstract}
직전 세 노트는 전부 \textbf{채점 규칙}을 다뤘고 효과는 한 번도 안 냈다.
이 노트는 규칙을 한 글자도 안 고치고 그 규칙이 뽑은 \textbf{A 등급 짝 둘}을 골라
효과와 위약을 끝까지 냈다.
\claimx{A 등급은 처치--결과 관계에 대해 \textbf{0 비트}를 나른다.}
{처치 열을 섞었을 때 절단값$\cdot$저군$\cdot$고군$\cdot$소수 쪽 중 \emph{하나라도} 달라지면 무너진다.
실측은 두 팔 다 \textbf{2400/2400 재현}이다.}
등급을 매기는 함수는 \textbf{인자에 결과가 없다} --- 값의 다중집합만 본다.
그래서 처치 열을 무작위로 섞고 \emph{똑같은 규칙}으로 다시 가르면
가짜 분할이 진짜 분할을 완전히 재현한다.
\claimx{그리고 효과 쪽은 위약 한복판에 있다.}
{구간이 0 을 배제하거나 $|\delta|$ 가 위약 95 분위를 넘으면 무너진다. 둘 다 아니었다.}
그래서 이 노트의 판정문은 \textbf{「효과가 없다」가 아니라 「A 등급이 효과를 잴 힘을 안 준다」}이고,
다음 사이클은 채점기가 아니라 \textbf{표본}으로 간다.
\end{abstract}

\section{왜 재는 쪽으로 갔나}

직전 노트는 「식별 가능(A) 짝이 몇 개인가」를 \textbf{세 번 다르게} 세었다.
프레시 티처가 마지막 판을 풀어 보니 조건 넷 중 둘은 상수, 하나는 표본 크기 검사로 퇴화,
남은 하나는 도메인 상수와 동어반복이었다.
\textbf{네 번째로 채점기를 고치면 A 는 또 다른 수가 되고, 그게 자료인지 코드인지는 또 다음 티처가 판정한다.}
이 되먹임을 끊는 방법은 채점을 더 잘하는 것이 아니라 \textbf{채점의 결과를 한 번 써 보는 것}이다.

사전등록 \texttt{docs/prereg\_903\_effect.md} 는 커밋 \texttt{2bd5532f4}
(2026-08-11T06:30:52+09:00)로 들어갔고 측정 산출물 스탬프는 06:38:14$\to$06:38:22 이다.
\texttt{stamp} 은 \textbf{실행 시작}에서 부른다.
러너 \texttt{runners/effect903.py} 는 채점기를 \textbf{import 만} 한다 ---
\texttt{median\_split}$\cdot$\texttt{sik1}$\cdot$\texttt{typed\_values}$\cdot$\texttt{d3\_backtrack}$\cdot$\texttt{resolve\_d3}$\cdot$\texttt{MIN\_SIDE}.
\textbf{등급을 다시 매기지 않았다.}

\section{먼저 원장에게 물었다}

이 저장소는 「없는 원천을 찾는」 단계는 갖고 있었지만
\textbf{이미 저장소 안에 있는 판정을 찾는 단계}는 없었다.
직전 사이클이 그 구멍에 빠졌다 --- 사다리가 겨누는 자가 이미 「결정에 안 번역된다」로
판정 나 있었는데 아무도 안 봤다.

그래서 재기 전에 원장(최상위 항목 \textbf{889} $\cdot$ 리프 \textbf{9,843})을 훑어
이 물음에 걸리는 \textbf{옛 판정 여덟}을 세었다. 둘만 옮긴다.

\begin{itemize}
\item 이 열은 \textbf{이미 챔피언 판의 공유 축}이고, 한 노트가 방향을 되돌려 \textbf{채택}까지 갔다.
\item 그 노트가 \textbf{팝업을 「도메인 안 부호가 음수인 곳」}에 넣어 뒀다.
\end{itemize}

\claimx{그래서 「팝업 $\delta$ 의 부호는 음수다」를 반증 가능한 예측으로 미리 박았고, 부호는 맞았다.}
{$\delta \geq 0$ 이면 무너진다. 실측 $\delta = -0.0244$.}
⚠ 다만 \textbf{자가 다르다} --- 옛 판정의 부호는 판 라벨과 눈금 맞춘 축에 대한 것이고
이 노트는 원천의 원값과 원 방문자 수를 쓴다. \textbf{같은 이름의 다른 자다.}
그리고 부호가 맞아도 크기는 위약 안이다(다음 절).

\section{짝을 재기 전에 고정했다 --- 그리고 지시서의 모순을 재기 전에 닫았다}

\begin{center}
\begin{tabular}{llrl}
\toprule
팔 & 처치 $\to$ 결과 & 짝 & 분모 딱지 \\
\midrule
ㄴ & 팝업 \texttt{entry\_friction} $\to$ 방문자 & 95 & \textbf{D1} 원천 레코드 380 \\
ㄱ & 영화 \texttt{개봉일} $\to$ 판 축의 $y$ & 406 & \textbf{D3} 판 채점 유보(챔피언 배선) \\
\bottomrule
\end{tabular}
\end{center}

\textbf{두 수를 잇지 않는다} --- D1 과 D3 는 다른 분모다.
팝업은 축 파일이 없어 D3 되짚기가 \textbf{원리상 불가능}하다. 「0」이 아니라 \textbf{「못 셌다」}이다.

🔴 처방서가 스스로 어긋나 있었다. ① 이 영화 \texttt{개봉일} 을 지정하면서
④ 는 「날짜형 아홉은 이번에 안 쓴다 --- 처치가 곧 결과 창의 원점이라」고 적는데,
\texttt{개봉일} 이 \textbf{바로 그 아홉 중 하나}다.
사전등록이 \textbf{재기 전에} 이렇게 닫았다: \textbf{재되 그것으로 판정하지 않는다} $\cdot$
나머지 여덟은 안 잰다 $\cdot$ 판정은 팝업 팔이 진다 $\cdot$
그리고 \textbf{그 교란을 주장이 아니라 측정으로 낸다.}

\claimx{창 절단은 실재한다 --- 늦게 개봉한 절반의 관측 창이 실제로 덜 찼다.}
{차가 0 이거나 부호가 반대면 「구조적 교란」주장이 이 자료에서 무너진다.
실측 차 $-0.3857$ 일차 $\cdot$ 검열 비율 3.5\% $\to$ 5.8\% $\cdot$ 관측일차와 $y$ 의 스피어만 $+0.319$.}

\section{배선을 먼저 부쉈다}

측정 전에 결함 넷을 \textbf{심어서} 게이트가 실제로 발화하는지 봤다. 넷 다 발화했고,
\textbf{음성 대조(정상 배선)도 넷 다 통과}했다.

\begin{center}
\begin{tabular}{lllc}
\toprule
심은 것 & 게이트 & 정상 $\to$ 심은 뒤 & 발화 \\
\midrule
분할 함수 파손 & 양쪽 $\geq 10$ & 저군 41 $\to$ 0 & ✓ \\
위약 셔플 무효화 & 값 가짓수 $\geq 2$ & 49 $\to$ 1 & ✓ \\
키 정규화 되돌리기 & 되짚기 조용한 0 & 406 $\to$ 0 & ✓ \\
결과 열 오타 & 짝 행 $\geq 20$ & 95 $\to$ 0 & ✓ \\
\bottomrule
\end{tabular}
\end{center}

마지막 하나가 이 저장소가 제일 무서워하는 꼴이다 --- \textbf{빈 것을 부정으로 읽는 길}.

\section{효과와 위약}

\begin{figure}[t]
\centering
\includegraphics[width=\linewidth]{figs/fig1.pdf}
\caption{(a) 두 팔 다 구간이 0 을 문다. 회색 띠는 무작위 중앙값 분할 위약의 $|\delta|$ 95 분위.
(b) 진짜 효과가 위약 분포 한복판에 있다 --- 위약이 진짜를 넘은 비율 0.840.
(c) \textbf{처치를 아무렇게나 섞어도 등급이 그대로다}: 위약 ② 의 분할 재현 2400/2400(두 팔)
$\cdot$ 위약 ③ 의 식1$=$예 2400/2400.}
\end{figure}

효과 크기는 Cliff's $\delta$(저군 대 고군 $\cdot$ 순위 기반), 구간은 규약 47 의 군집 부트스트랩
\textbf{BCa} 다. 두 팔 다 \textbf{폴백 없이 BCa} 가 나왔다.

\begin{center}
\begin{tabular}{lrrlr}
\toprule
팔 & $\delta$ & BCa 95\% & 판정 & 위약 ② $|\delta|_{95}$ \\
\midrule
팝업 (판정) & $-0.0244$ & $[-0.2740, +0.2013]$ & 판정 불능 & 0.2372 \\
영화 (판정 안 함) & $-0.0487$ & $[-0.1629, +0.0633]$ & 판정 불능 & 0.1115 \\
\bottomrule
\end{tabular}
\end{center}

군집은 팝업이 브랜드 열쇠(95 행 $\to$ 83 군집 $\cdot$ 병합 12),
영화가 프랜차이즈 열쇠(406 행 $\to$ 404 군집 $\cdot$ 병합 2)다.
\textbf{병합이 적으면 구간 폭이 과소 방향}이고 그 병기는 자가 스스로 붙인다.

위약은 셋이다: 결과를 섞는 것 $\cdot$ \textbf{처치 열을 섞고 같은 규칙으로 다시 가르는 것} $\cdot$
결측까지 포함해 원천 전량을 섞는 것. 씨앗 12 $\times$ 반복 200 $=$ 2,400.
씨앗별 평균의 표준편차는 네 조합에서 0.0038$\sim$0.0088 로 씨앗 잡음이 결론을 흔들 자리에 없다.

\subsection{그래서 무엇이 0 비트인가}

\claimx{등급을 매기는 함수는 결과를 한 번도 안 본다.}
{$\texttt{sik1}$ 이나 $\texttt{median\_split}$ 의 인자에 결과가 하나라도 있으면 이 문장은 거짓이다.}
그 결과 처치 열을 어떻게 섞어도 절단값과 두 군의 크기가 그대로다 ---
팝업 2400/2400 $\cdot$ 영화 2400/2400 $\cdot$ 결측까지 섞은 위약 ③ 의 식1$=$예 2400/2400.

\textbf{이것은 항등식이고, 그 사실을 알고도 쟀다.}
직전 티처가 잡은 병은 항등식을 \emph{증거로 판 것}이었지 항등식을 \emph{잰 것}이 아니다.
여기서는 \textbf{항등식이라는 것 자체가 결론}이다 ---
「A 21」은 「자료가 개입을 식별한다」가 아니라
\textbf{「그 열의 값이 골고루 퍼져 있다」}는 말이다.
티처가 서로 다른 다섯 각도에서 가리킨 것이 같은 하나였다.

한 노트가 문턱을 두고 \textbf{「위약이 아무것도 안 했다 --- 문턱은 위약으로 잴 수 있는 물건이
아니었다」}고 적은 적이 있다. \textbf{A 등급도 그렇다.}

\section{그러면 자는 얼마나 힘이 없나}

\begin{figure}[t]
\centering
\includegraphics[width=\linewidth]{figs/fig2.pdf}
\caption{(a) 이 설계가 볼 수 있는 최소 효과(MDE) 대 실측 $|\delta|$.
팝업은 실측이 MDE 의 \textbf{7.2\%}, 영화는 30.4\%.
(b) 창 절단을 주장이 아니라 측정으로 --- 늦게 개봉한 절반의 관측 창이 실제로 덜 찼다.}
\end{figure}

「힘을 안 준다」를 수로 바꾸면 이렇다(사후 서술 $\cdot$ 판정 미사용):
팝업 팔이 볼 수 있는 최소 효과는 $|\delta| \approx 0.339$ 인데 실측은 0.0244 다.
\textbf{자가 못 보는 크기의 것을 재려 하고 있었다.}

\claimx{그래서 판정문은 이 문장이다 --- 「효과가 없다」가 아니라 「A 등급이 효과를 잴 힘을 안 준다」.}
{구간이 0 을 배제했거나 $|\delta|$ 가 위약 95 분위를 넘었으면 이 문장을 못 쓴다. 둘 다 거짓이었다.}
\textbf{다음 사이클은 채점기가 아니라 표본으로 간다.}

\section{자가 적발 둘}

\textbf{하나.} 사전등록의 예측 P5(「두 위약 분포는 구별 안 된다」)의 판정 규칙이
\textbf{거의 동전 던지기}였다 --- 두 분포가 같은 법이면 한쪽의 95 분위가 다른 쪽의 95 분위를
넘을 확률이 대략 절반이다. 실제로 팝업은 「반증」, 영화는 「확인」이 나왔고
팝업의 두 위약 표준편차는 0.1194 와 0.1210 이다. \textbf{규칙을 안 고치고 결과를 그대로 실었다} ---
고치면 그게 사후 규칙 변경이다.

\textbf{둘.} 이 사이클이 낸 구간은 넷이고 그중 0 을 안 문 것은 \textbf{하나}인데,
그 하나가 사전등록이 \emph{측정 전에} 「부차 $\cdot$ 판정 미사용」으로 못 박은 팔이다
(짝 35 $\cdot$ 매출을 보고한 팝업만 남은 것이라 보고 선택이 그대로 살아 있다).
\textbf{이 노트의 결론에 안 넣는다} --- 넣으면 그게 사후 부분집합이다.
다중성 보정은 \textbf{안 했다}(안 한 것이지 필요 없는 것이 아니다):
넷 중 하나가 0 밖인 것은 우연으로도 흔하다.

\section{못 잰 것 --- 「없다」가 아니다}

\begin{itemize}
\item 팝업의 D3 되짚기 --- 축 파일이 없어 \textbf{원리상 불가능}하다
\item 날짜형 나머지 여덟 짝 --- 사전등록대로 \textbf{안 쟀다}
\item 교란 통제 --- 어느 팔도 공변량을 하나도 안 넣었다. $\delta$ 는 \textbf{조건 없는 두 군 차}다
\item 비단조 --- 한 옛 노트가 「규모가 부호를 바꾼다」고 적어 뒀는데
      \textbf{중앙값 한 번의 분할로는 그것을 못 본다}
\end{itemize}

\section{남는 것}

이 저장소는 「A 등급 짝이 있다」를 세 번 다르게 세었고, 이 노트가 처음으로
\textbf{그 등급이 무엇을 가능하게 하는지}를 봤다. 답은 \textbf{아무것도}다 ---
등급은 결과를 안 보고 매겨지므로 위약이 100\% 재현하고,
효과 쪽은 자가 볼 수 있는 크기의 7\% 다.
\textbf{다음에 늘려야 하는 것은 조건의 수가 아니라 짝의 수다.}

\end{document}
